\subsection{The normalized Fourier bound at every actual head}
\label{actual-kernel-fourier}

We now obtain the original one-coordinate Fourier estimate from
the actual coefficient maps.  The point normalization, terminal
density and old-word rational factors are retained throughout.
This result concerns a supplied original place and top lift; it
does not construct that arithmetic place or its integral lattices.

Fix the original data $P,p$ with $k\ge2$, its normalized tower,
compact types, base frame, and top lift $\Theta$.  Retain the
original regular-radial top-lift identity and
$\mathrm{InDDplus2}(C_b)$ for a fixed $C_b$.
Let $F$ be a smooth compact coefficient with positive Gram cutoff.
The word operators, root quotient measures, and forms are the
original ones.  Put
\[
 q_0(Y)=-TT'\prod_{i=1}^k
       (Y+\tfrac12+u_i)(Y+\tfrac12+u'_i),\qquad
 S_w=\prod_{(\varepsilon,a)\in w}n(a),
\]
where $n$ is strictly positive at all parameters used below.
Here $n(a)=|P_{\mathrm{eig}}(a)|$ in the original eigenfunction
application; it is not the dimension of a representation.
For a head $h$ of size $j\le r$, write
\[
 a_h=n\mathrm{Pt}(k,r,TT',h)>0,\qquad
 x_{w,h}=\bigl(S_h\mathsf V_{j,w,e}F\bigr)_e,
 \qquad \beta_{j,w}=\mathrm{formLev}_{j,w}.
\]
Thus $\mathrm{RR}_{j,w}(h,h')=
\beta_{j,w}(x_{w,h},x_{w,h'})$.
Denote the original normalized kernel and its diagonal by
\[
 N_{i,w}(a,b)=\frac{\sqrt{n\mathrm{Pt}(a)}
                          \sqrt{n\mathrm{Pt}(b)}}{S_w}
                       \mathrm{RR}_{i,w}(a,b),\qquad
 E_{i,w}(a)=\operatorname{Re}N_{i,w}(a,a).
\]
These are exactly the source's $\mathrm{NP}$ and $\mathrm{Ew}$
\cite{kernel-fourier-source}.

Use the already constructed finite grids
$\mathcal G_0,\mathcal G_1,\ldots,\mathcal G_N$ in
$[\alpha,\beta]$ contained in the original cut, with
\[
 N=20(r+1),\qquad I=\beta-\alpha\ge6N^3+6N,
 \qquad \mathcal Y=\bigcup_{a=0}^N\mathcal G_a.
\]
Every $y,y+1$ belongs to the original cut.  The short grid
has $N$ points.  Different long grids, and a long and a
short grid, are separated by at least $N^2+1$ after either
zero or one shift; distinct short nodes have the lower bound
$N^2/8$.  Each long grid contains $2j+1$ interpolation
nodes whose sum of absolute leading-coefficient weights is
at most $(28N/I)^{2j}$, for $1\le j\le r$.
These are the original grid ranges; no boundedness of $r$
is assumed \cite{kernel-fourier-established}.

Set
\[
 A_0=\max_{y\in\mathcal G_0}
       \sqrt{\frac{n(y)n(y+1)}{q_0(y)}},\qquad
 A_1=\max_{y\in\mathcal Y}
       \sqrt{\frac{n(y)n(y+1)}{q_0(y)}}.
\]
The original cut gives $q_0(y)>0$, so these are finite and
positive.  Define
\begin{equation}\label{kernel-fourier-Lambda}
 \Lambda_j=
 \begin{cases}
 e^{10/49}A_0,&j=0,\\
 e^{10/49}A_1(28N/I)^{2j},&1\le j\le r,
 \end{cases}
 \qquad \Lambda=\max_{0\le j\le r}\Lambda_j.
\end{equation}
One may enlarge $\Lambda$, in particular to $\max(1,\Lambda)$
when the downstream statement requires it.

\begin{theorem}\label{kernel-fourier-actual}
For every integer $0\le K\le N$, every $j\le r$, every
relevant head $h$ of size $j$, and every original cut word
$w$ of length $|w|+K\le N$ whose letters are
$(+,y)$ or $(-,y+1)$ with $y\in\mathcal Y$, the following holds.
For every $s\in\mathbb C$ there is a word $w'=L\mathbin{\|}w$
on that same grid, satisfying the same original cut and
separation, with $|L|\le K$ and $|w'|\le N$, such that
\begin{equation}\label{kernel-fourier-bound}
 \left|\int_{\mathbb C} e^{sc-\bar s\bar c}
 N_{j+1,w}\bigl(h\mathbin{\|}(1,c),h\mathbin{\|}(1,c)\bigr)
 \,dc\right|
 \le m(\Lambda,K,|s|)E_{j,w'}(h),
\end{equation}
where $m(\Lambda,K,0)=1$ and
$m(\Lambda,K,t)=\min(1,(\Lambda/t)^K)$ for $t>0$.
The density has a continuous integrable representative agreeing
with it at every $c\ne0$.
All original $E_{i,w}(a)$ are nonnegative, including terminal
and invalid heads.  Thus this supplies the original
\texttt{KernelFourierIntegralBound},
\texttt{HeadContinuityAtZero} and \texttt{PosAt} on these grids.
\end{theorem}

\begin{proof}
First identify the measure and the head normalization.
The original recursion for $n\mathrm{Pt}$ gives, for $j<r$,
\[
 n\mathrm{Pt}\bigl(h\mathbin{\|}(1,c)\bigr)
       =a_h|c|^{2(2(r-j)+k-1)}.
\]
At $j=r$, the factor is instead
$(|c|^2-TT')^{k-1}$ on $|c|^2>TT'$.
The original orbital vectors vanish outside this positive
domain, including its boundary.  Therefore negative values
of the formal power outside the domain do not enter a
square-root simplification.  Both cases give the same
normalized integral identity
\begin{equation}\label{kernel-fourier-pairing}
 \widehat H_{w,h}(s)
   =\frac{a_h}{S_w}
      \beta_{j,w}(U_sx_{w,h},x_{w,h}),
\end{equation}
where $\widehat H$ is the integral in
\eqref{kernel-fourier-bound}, and $U_s$ is the actual
phase together with its prescribed triangular tensor action.

For $j<r$, this is the proved partial-head transfer and
one-coordinate radial Fourier formula.  For $j=r$, the
proved terminal transfer identifies the actual terminal
vectors and all their type weights; integration over the
original normalized terminal group pushes forward to
$(|c|^2-TT')^{k-1}dc$ on the positive domain.  Its compact
support away from the boundary gives a continuous compactly
supported density.  Below the terminal level, the proved
Schwartz-marginal theorem supplies a continuous integrable
extension, identified pointwise on $c\ne0$ by the actual
compact-family regularity.  Multiplication by the fixed
positive scalar $a_h/S_w$ proves the stated continuity.
The value at zero has measure zero; every nonzero point,
including every relevant rational point, is covered by
the actual identity.  No ratio of Haar constants occurs
in \eqref{kernel-fourier-pairing}.

The forms in \eqref{kernel-fourier-pairing} are positive
semidefinite on the actual compact coefficient class.
The proved wave identity preserves them and that class.
Cauchy--Schwarz therefore gives
\begin{equation}\label{kernel-fourier-unweighted}
                |\widehat H_{w,h}(s)|\le E_{j,w}(h).
\end{equation}
The same positivity gives nonnegative $E$ below the terminal
level.  At the terminal level its definition is a sum of
positive type weights times Hilbert--Schmidt squares; invalid
heads have value zero.  This proves the full stated positivity.

We must choose extension grids for \emph{every} word in
the theorem.  No position condition on the old grid blocks
is being assumed.  Assign to each old letter its underlying
node $a$ for $(+,a)$ and $a-1$ for $(-,a)$.  The grids
are disjoint.  At most $|w|$ long blocks are occupied by
old underlying nodes; hence at least $K$ long blocks are
unoccupied.  If $j\ge1$, choose any $K$ distinct such
blocks and take their degree-$2j$ interpolation nodes.
For $j=0$, at most $|w|$ points of the short grid are
occupied; choose $K$ distinct unused short nodes, each
as a one-point degree-zero interpolation grid.  Every
choice of one node in each selected grid then produces
both admissible extensions
\[
 w_+(\boldsymbol y)=((+,y_1),\ldots,(+,y_K),w),\qquad
 w_-(\boldsymbol y)=((-,y_1+1),\ldots,(-,y_K+1),w).
\]
Their length is $|w|+K$, their parameters remain in the
original cut, and every new--old or new--new parameter
distance is greater than two.  This argument also covers
old words with several letters in a single long block.

For each selected node the exact modified determinant is
\[
 q_w(y)=q_0(y)
  \prod_{(+,a)\in w}\frac{y-a+1}{y-a}
  \prod_{(-,a)\in w}\frac{a-y}{a-y-1}.
\]
Its denominators have modulus at least $D=N^2/8$.
The already proved logarithmic estimate gives
$|\log(q_w(y)/q_0(y))|\le |w|/(D-1)\le20/49$.
Consequently
\begin{equation}\label{kernel-fourier-q-ratio}
 \sqrt{\frac{n(y)n(y+1)}{q_w(y)}}
       \le e^{10/49}A_0\quad(j=0),\qquad
       \le e^{10/49}A_1\quad(j>0).
\end{equation}
Every rational correction from the old word is included here.

Write $c_{a,i}$ for the Lagrange leading-coefficient weights,
$C_a=\sum_i|c_{a,i}|$, and
$\theta_{\boldsymbol i}=\prod_a|c_{a,i_a}|/\prod_a C_a$.
These are nonnegative and sum to one.  The proved actual
Capelli leading-coefficient identity, with degree $2j$ at
each insertion, gives exactly
\[
 s^K\beta_{j,w}(U_sx_{w,h},x_{w,h})
  =\sum_{\boldsymbol i}\left(\prod_a c_{a,i_a}\right)
       \mathcal B_w\bigl(U_sx_{w_+(\boldsymbol y),h},
                              x_{w_-(\boldsymbol y),h}\bigr).
\]
Lists are processed tail first.  The negative insertion is
at $y+1$, and the mixed contraction is
$\beta(x_1,y_0)-\beta(x_0,y_1)$.  Thus the displayed
coefficient is $s^K$, with no conjugation or sign loss.
All polynomial coefficients lie in the actual compact
coefficient class: the slice theorem and finite Vandermonde
inversion give this directly, without interchanging an
unjustified derivative and integral.

Apply the actual mixed determinant inequality to each
summand and remove $U_s$ from its positive diagonal by
the actual wave isometry.  The word normalizers obey the
exact identities
\[
 S_{w_+}=S_w\prod_a n(y_a),\qquad
 S_{w_-}=S_w\prod_a n(y_a+1).
\]
After multiplication by $a_h/S_w$, the factors $a_h$ and
$S_w$ cancel against these normalizers.  Using
\eqref{kernel-fourier-q-ratio}, and then the triangle
inequality, yields
\begin{equation}\label{kernel-fourier-weighted}
 |s|^K|\widehat H_{w,h}(s)|
  \le \Lambda_j^K\sum_{\boldsymbol i}\theta_{\boldsymbol i}
       \sqrt{E_{j,w_+(\boldsymbol y)}(h)
                     E_{j,w_-(\boldsymbol y)}(h)}.
\end{equation}
This is an estimate for the original normalized integral,
not just an abstract form.

Let $\mathcal E$ consist of $w$ and the finitely many
positive and negative extensions just constructed.
Choose $w'\in\mathcal E$ maximizing $E_{j,w'}(h)$, and
write this nonnegative maximum as $M$.  Each geometric
mean in \eqref{kernel-fourier-weighted} is at most $M$,
while \eqref{kernel-fourier-unweighted} is at most $M$.
For $s\ne0$ take the stronger of these two bounds and
use $\Lambda_j\le\Lambda$.  At $s=0$ use
\eqref{kernel-fourier-unweighted}.  When $K=0$ the
extension set is just $\{w\}$ and the same conclusion
holds with every empty product equal to one.  Every
element of $\mathcal E$ satisfies the original
$\mathrm{IsExt}(N,K,w,\cdot)$ contract.  This proves
\eqref{kernel-fourier-bound} with all its quantifiers.
\end{proof}

\paragraph{Arithmetic scale and remaining inputs.}
For the prescribed eigenpolynomial take $n(y)=|P_{\mathrm{eig}}(y)|$.
Let $m=2r+2$, retain the original $g,\kappa>0$, and write
$t_v=\sqrt{T_{\epsilon,v}T_{\epsilon',v}/(TT')}$.
The established boundary-polynomial theorem gives
\[
 A_0\le g^{-m}\kappa^{m-1}\,80t_v\sqrt{k}
                       \sqrt{N^3/\Delta},\qquad
 A_1\le g^{-m}\kappa^{m-1}\,80t_v\sqrt{k}.
\]
Its hypotheses include the original binary spectral
construction and a differing index, together with the
stated interval condition.  Inserting these bounds in
\eqref{kernel-fourier-Lambda} preserves every arithmetic
parameter and every power of $g,\kappa,I,N$.
For the original character
$\exp(4\pi i\operatorname{Re}(\ell c)/g)$ and the original
variable change $z=\lambda c$, the wave parameter is
$s=2\pi i\ell/(g\lambda)$; hence its frequency scale is
exactly $(g\lambda/(2\pi))\Lambda_j$.  The Jacobian of the
same change belongs to the original normalized integral
and is not absorbed in this frequency scale.

This closes the actual one-coordinate Fourier supplier on
compact coefficients.  Passing to the prescribed decaying
eigenfunctions, establishing the remaining block-growth
comparison, constructing the original uniformly controlled
top lift, and identifying the arithmetic source and target
Gram operators remain separate obligations.  None is
replaced by a bound on this finite extension maximum.
